\documentclass[11pt]{amsart}
\usepackage[T1]{fontenc}
\usepackage{lmodern,microtype,amsmath,amssymb,amsthm,mathtools,enumitem,booktabs}
\usepackage[margin=1.1in]{geometry}
\usepackage[hidelinks]{hyperref}
\usepackage{xurl}

\newtheorem{theorem}{Theorem}[section]
\newtheorem{lemma}[theorem]{Lemma}
\newtheorem{corollary}[theorem]{Corollary}
\newtheorem{proposition}[theorem]{Proposition}
\theoremstyle{remark}
\newtheorem{remark}[theorem]{Remark}
\numberwithin{equation}{section}

\newcommand{\dd}{\,\mathrm{d}}

\hypersetup{pdftitle={The Lyapunov exponent of a randomly switched pair of stable planar systems is unimodal in the switching rate},pdfauthor={Yitzchak Shmalo}}

\title[Unimodality of a switched Lyapunov exponent]{The Lyapunov exponent of a randomly switched pair of stable planar systems is unimodal in the switching rate}
\author{Yitzchak Shmalo}
\address{Einstein Institute of Mathematics, The Hebrew University of Jerusalem, Jerusalem, Israel}
\email{yitzchak.shmalo@gmail.com}
\date{5 November 2026}
\subjclass[2020]{60J25, 37H15, 93E15, 65G30}
\keywords{piecewise deterministic Markov process, switched linear system, Lyapunov exponent, interval arithmetic}

\begin{document}
\begin{abstract}
Let $X_t$ solve $\dot X_t=A_{I_t}X_t$ in the plane, where $A_0=\begin{psmallmatrix}-\alpha&1\\0&-\alpha\end{psmallmatrix}$, $A_1=\begin{psmallmatrix}-\alpha&0\\-1&-\alpha\end{psmallmatrix}$ and $I_t$ switches between $0$ and $1$ at rate $r$. Both matrices are Hurwitz, yet for small $\alpha$ the system blows up for an intermediate range of rates. The top Lyapunov exponent has the form $G(r)-\alpha$, and Malrieu conjectured that $G$ increases and then decreases, with a single maximum of height about $0.2$, and tends to zero at both ends. We prove this. The maximum is attained at a single $r_*\in(0.226,0.235)$, $G(r)=r\log(1/r)\,(1+O(1/\log(1/r)))$ as $r\to0$, and $G(r)\le 1/(2r)+(2/\pi)K_1(r)$ for all $r$. Consequently the set of rates for which the system is unstable is an open interval when $\alpha<G(r_*)$ and is empty otherwise. The proof writes $G$ as a ratio of two one-dimensional integrals of an explicit kernel, settles $r\le 0.01$ and $r\ge5$ with elementary and Bessel-function bounds, and covers $[0.01,100]$ by interval arithmetic.
\end{abstract}
\maketitle

\section{Introduction}\label{sec:intro}

Fix $\alpha>0$ and $r>0$ and consider the Markov process $(X_t,I_t)$ on $\mathbb R^2\times\{0,1\}$ with generator
\begin{equation}\label{eq:gen}
Lf(x,i)=(A_ix)\cdot\nabla_xf(x,i)+r\bigl(f(x,1-i)-f(x,i)\bigr),\qquad
A_0=\begin{pmatrix}-\alpha&1\\0&-\alpha\end{pmatrix},\quad A_1=\begin{pmatrix}-\alpha&0\\-1&-\alpha\end{pmatrix}.
\end{equation}
Every convex combination $pA_1+(1-p)A_0$ has eigenvalues $-\alpha\pm i\sqrt{p(1-p)}$, so each frozen system is exponentially stable, but Lawley, Mattingly and Reed \cite{LMR} showed that for suitable $\alpha$ the switched system is unstable for an intermediate range of rates. Write $U_t=X_t/|X_t|$. The pair $(U_t,I_t)$ has a unique invariant law $\mu_r$, and almost surely
\[
\lim_{t\to\infty}\frac1t\log|X_t|=\sum_{i}\int\langle A_iu,u\rangle\,\mu_r(\mathrm du,i)=G(r)-\alpha,
\]
where $G$ does not depend on $\alpha$ and has an explicit integral formula \cite[Theorem 4.1]{Ma}. Malrieu \cite[Conjecture 4.3]{Ma} conjectured from numerical experiments that $G$ is unimodal: there is $r_c$ with $G'>0$ on $(0,r_c)$ and $G'<0$ on $(r_c,\infty)$, $G(r_c)\approx0.2$, and $G(r)\to0$ as $r\to0$ and as $r\to\infty$. He noted that this would make the instability set $\{r>0:G(r)>\alpha\}$ an interval for $\alpha<G(r_c)$ and empty for $\alpha>G(r_c)$, and asked for a proof \cite[Open question 2]{Ma}. For a family of rotation-type pairs, Lagasquie \cite{La} proves that the exponent tends to its frozen value at both ends of the rate axis.

\begin{theorem}\label{thm:main}
$G$ is continuously differentiable on $(0,\infty)$, and there is $r_*\in(0.226,0.235)$ such that $G'>0$ on $(0,r_*)$ and $G'<0$ on $(r_*,\infty)$. Moreover
\[
\frac{\log(1/r)-\gamma}{1+4r^2(1+2K_0(2r))}\ \le\ \frac{G(r)}{r}\ \le\ \log\frac1r+2K_0(2-2r^2)\qquad(0<r<1),
\]
where $\gamma$ is Euler's constant, and $G(r)\le \frac1{2r}+\frac2\pi K_1(r)$ for all $r>0$. In particular $G(r)\to0$ at both ends.
\end{theorem}

Numerically $G(r_*)=0.1985$. The two-sided bound gives $G(r)=r\log(1/r)\,(1+O(1/\log(1/r)))$; the quadrature values are consistent with $G(r)=r\log(1/r)-\gamma r+O(r^3\log^2(1/r))$.

\begin{corollary}\label{cor:interval}
Let $\alpha>0$. If $\alpha<G(r_*)$, there are $r_1<r_*<r_2$ such that the Lyapunov exponent $G(r)-\alpha$ is positive, and hence $|X_t|\to\infty$ almost surely, exactly for $r\in(r_1,r_2)$, and negative, and hence $|X_t|\to0$ almost surely, for $r\notin[r_1,r_2]$. If $\alpha\ge G(r_*)$, the exponent is negative for every $r\ne r_*$.
\end{corollary}

This follows from the theorem, since a continuous function that increases strictly, then decreases strictly, and tends to $0$ at both ends has interval superlevel sets.

\begin{remark}\label{rem:scale}
In \cite{Ma} the maximum is placed at $r_c\approx4.6$, read from Figure~1 there. For the system \eqref{eq:gen} the maximum is at $r_*\approx0.23$. A direct simulation of \eqref{eq:gen} confirms this: integrating the two shear flows exactly over independent exponential holding times of rate $r$, with $2\times10^6$ switches, gives $\frac1T\log|X_T|+\alpha=0.1985$, $0.1081$, $0.0272$, $0.0065$ at $r=0.23$, $1$, $4.6$, $20$, against $G=0.1985$, $0.1083$, $0.0269$, $0.0062$ from the formula of Section~\ref{sec:reduction}. The shape, the height and the two limits are as conjectured.
\end{remark}

The proof has three parts. Section~\ref{sec:reduction} rewrites $G$ in the coordinate $x=\cot\theta$, uses a symmetry exchanging the two regimes to reduce to $x>0$, and arrives at $G=N_+/D_+$ with
\begin{equation}\label{eq:ND}
N_+(r)=\int_0^\infty\Bigl(\frac{x}{1+x^2}-\frac{J(x,r)}x\Bigr)\dd x,\qquad D_+(r)=\int_0^\infty\Bigl(1-J(x,r)+\frac{J(x,r)}{x^2}\Bigr)\dd x,
\end{equation}
\begin{equation}\label{eq:J}
J(x,r)=rx\int_0^1e^{-r\varphi(x,s)}\dd s,\qquad \varphi(x,s)=xs+\frac{s}{x(1-s)}.
\end{equation}
Since $G'=(N_+'D_+-N_+D_+')/D_+^2$, the sign of $G'$ is the sign of $H_+=N_+'D_+-N_+D_+'$. Section~\ref{sec:small} proves $H_+>0$ for $r\le0.01$, Section~\ref{sec:large} proves $H_+<0$ for $r\ge5$, and Section~\ref{sec:middle} covers $[0.01,100]$ by validated numerics.

\section{Reduction to a ratio of integrals}\label{sec:reduction}

Since $A_i=-\alpha I+N_i$ with $N_0=\begin{psmallmatrix}0&1\\0&0\end{psmallmatrix}$ and $N_1=\begin{psmallmatrix}0&0\\-1&0\end{psmallmatrix}$, the direction $U_t$ moves by the shears $N_i$ alone, and $\langle A_iu,u\rangle=-\alpha+\langle N_iu,u\rangle$. The law $\mu_r$ is invariant under $u\mapsto-u$, so we may work on the projective line with the coordinate $x=u_1/u_2\in\mathbb R\cup\{\infty\}$. In regime $0$ the flow $\dot u=N_0u$ gives $\dot x=1$, in regime $1$ the flow $\dot u=N_1u$ gives $\dot x=x^2$, and
\[
\langle N_0u,u\rangle=\frac{x}{1+x^2},\qquad \langle N_1u,u\rangle=-\frac{x}{1+x^2}.
\]
Let $f_0,f_1$ be the densities of the invariant law of $(x,I)$, up to a common factor. They satisfy $-f_0'+r(f_1-f_0)=0$ and $-(x^2f_1)'+r(f_0-f_1)=0$, so the flux $f_0+x^2f_1$ is constant and may be taken equal to $1$. Then
\begin{equation}\label{eq:G0}
G(r)=\frac{\int_{\mathbb R}(f_0-f_1)\frac{x}{1+x^2}\dd x}{\int_{\mathbb R}(f_0+f_1)\dd x}.
\end{equation}

\emph{Symmetry.} The map $\Phi(x,i)=(-1/x,1-i)$ conjugates the dynamics to itself: if $\dot x=1$ then $y=-1/x$ has $\dot y=y^2$, and if $\dot x=x^2$ then $\dot y=1$; the switching rates are equal. By uniqueness, $\Phi$ preserves the invariant law, that is $f_1(y)=f_0(-1/y)/y^2$. The growth rate is also preserved, since $-y/(1+y^2)=x/(1+x^2)$ for $y=-1/x$. The half-line $\{x>0\}$ in regime $i$ is carried onto $\{x<0\}$ in regime $1-i$, so both integrals in \eqref{eq:G0} are twice their restrictions to $x>0$.

\emph{The density on $x>0$.} With $f_1=(1-f_0)/x^2$ the first equation reads $f_0'+r(1+x^{-2})f_0=r/x^2$. Its solutions on $(0,\infty)$ differ by multiples of $e^{-r(x-1/x)}$, which is not integrable at $0$, so
\[
f_0(x)=\int_0^x\frac r{y^2}\,e^{-r[(x-y)+(1/y-1/x)]}\dd y=\frac rx\int_0^1\frac{e^{-r\varphi(x,s)}}{(1-s)^2}\dd s,
\]
after the substitution $y=x(1-s)$. Since $\partial_s\varphi=x+1/(x(1-s)^2)$ and $\varphi$ increases from $0$ to $\infty$, integrating $r\,\partial_s\varphi\,e^{-r\varphi}$ over $s\in(0,1)$ gives $1=J+f_0$, that is $f_0=1-J$ with $J$ as in \eqref{eq:J}. Put $b=1+x^2$. Then $(f_0-f_1)\frac x{b}=\frac xb-\frac Jx$ and $f_0+f_1=1-J+J/x^2$, which with \eqref{eq:G0} gives $G=N_+/D_+$ with \eqref{eq:ND}.

\begin{lemma}\label{lem:J}
For $x,r>0$:
\begin{enumerate}[label=(\alph*)]
\item $1-e^{-rx}>J>0$ and $J\le x^2/b$;
\item $\partial_rJ=r\int_0^1\frac{s^2}{(1-s)^2}e^{-r\varphi}\dd s>0$.
\end{enumerate}
\end{lemma}

\begin{proof}
(a) $\varphi>xs$ gives $J<rx\int_0^1e^{-rxs}\dd s=1-e^{-rx}$. Also $\varphi-sb/x=s^2/(x(1-s))\ge0$, so $J\le rx\int_0^\infty e^{-rsb/x}\dd s=x^2/b$. (b) $\partial_rJ=x\int_0^1(1-r\varphi)e^{-r\varphi}\dd s$, and
\[
(1-r\varphi)e^{-r\varphi}=\partial_s\bigl(se^{-r\varphi}\bigr)+r\,(s\partial_s\varphi-\varphi)\,e^{-r\varphi},\qquad s\partial_s\varphi-\varphi=\frac{s^2}{x(1-s)^2}.
\]
The term $se^{-r\varphi}$ vanishes at $s=0$ and as $s\to1$.
\end{proof}

The integrands of \eqref{eq:ND} are nonnegative and integrable by Lemma~\ref{lem:J}(a) and Lemma~\ref{lem:large}(a) below. For $r\in[r_1,r_2]\subset(0,\infty)$, Lemma~\ref{lem:J}(b) gives $0<\partial_rJ(x,r)\le(r_2/r_1)\,\partial_rJ(x,r_1)$, and $\int_0^\infty(1+x^{-2})\,\partial_rJ(x,r_1)\dd x=r_1^{-2}$ by \eqref{eq:e3id} below. This dominates the $r$-derivatives $-\partial_rJ/x$ and $\partial_rJ\,(x^{-2}-1)$ of the integrands, so $N_+$ and $D_+$ are $C^1$ and may be differentiated under the integral sign. Hence $G\in C^1(0,\infty)$ and $\operatorname{sign}G'=\operatorname{sign}H_+$, with
\begin{equation}\label{eq:H}
H_+=N_+'D_+-N_+D_+',\qquad N_+'=-\int_0^\infty\frac{\partial_rJ}{x}\dd x,\qquad D_+'=\int_0^\infty\partial_rJ\Bigl(\frac1{x^2}-1\Bigr)\dd x.
\end{equation}

\section{Small rates}\label{sec:small}

The $x$-integrals against $e^{-r\varphi}=e^{-rs(x+a/x)}$, $a=1/(1-s)$, are Bessel functions:
\begin{equation}\label{eq:bessel}
\int_0^\infty x^{\nu-1}e^{-\beta(x+a/x)}\dd x=2a^{\nu/2}K_\nu(2\beta\sqrt a)\qquad(\beta,a>0),
\end{equation}
here with $\beta=rs$ and $2\beta\sqrt a=2rs/\sqrt{1-s}$. Put $u=s/\sqrt{1-s}$, which increases from $0$ to $\infty$, with $\dd u=\frac{2-s}{2(1-s)^{3/2}}\dd s$, and write $w=1-s$, so that $u^2w=(1-w)^2$, $0<w\le1$, $w\le u^{-2}$ and $\sqrt w=(1-w)/u$. We use the Mellin formula $\int_0^\infty v^{\mu-1}K_\nu(v)\dd v=2^{\mu-2}\Gamma(\frac{\mu-\nu}2)\Gamma(\frac{\mu+\nu}2)$, the bound $vK_1(v)\le1$, and $K_0'=-K_1$; the last two give $K_0(v)\le K_0(y)+\log(y/v)$ for $0<v\le y$, hence $\int_0^yK_0\le y(1+K_0(y))$.

Let $L=\log(1/r)$ and define $e_1,\dots,e_4$ by
\[
rD=2+e_1,\quad N=2L-2\gamma+e_2,\quad r^2D'=-2+e_3,\quad rN'=-2+e_4,
\]
where $N=2N_+$ and $D=2D_+$.

\begin{lemma}\label{lem:e}
For every $r>0$, $e_1,e_2,e_3,e_4>0$, and
\[
e_1\le8r^2\bigl(1+2K_0(2r)\bigr),\qquad e_3\le4r^2\bigl(1+2K_0(2\sqrt2\,r)\bigr),\qquad N\le2L+4K_0(2-2r^2)\quad(r<1).
\]
\end{lemma}

\begin{proof}
\emph{$e_1$.} By Fubini and \eqref{eq:bessel} with $\nu=0$, $\int_0^\infty(1-J)\dd x=\int_0^1\frac{2r}{(1-s)^2}K_0\dd s$ and $\int_0^\infty Jx^{-2}\dd x=\int_0^12rK_0\dd s$, the argument of $K_0$ being $2rs/\sqrt{1-s}=2ru$. In the variable $u$,
\[
D_+=4r\int_0^\infty u\,\frac{1+w^2}{1-w^2}\,K_0(2ru)\dd u .
\]
Since $4r\int_0^\infty uK_0(2ru)\dd u=1/r$, we get $e_1=2rD_+-2=16r^2\int_0^\infty\frac{uw^2}{1-w^2}K_0(2ru)\dd u>0$. Using $1-w=u\sqrt w$, $\frac{uw^2}{1-w^2}=\frac{w^{3/2}}{1+w}\le\min(\frac12,u^{-3})$, so
\[
e_1\le16r^2\Bigl(\frac1{4r}\int_0^{2r}K_0+\frac12K_0(2r)\Bigr)\le8r^2\bigl(1+2K_0(2r)\bigr).
\]
\emph{$e_2$.} By Lemma~\ref{lem:J}(a), $N_+>\int_0^\infty\bigl(e^{-rx}-\frac1{1+x^2}\bigr)\frac{\dd x}x=L-\gamma$. For the upper bound take $\delta=r^2$; on $s\le1-\delta$, $\varphi\le s(x+1/(x\delta))$, so
\[
J\ge\frac{\delta x^2}{1+\delta x^2}\Bigl(1-e^{-r(1-\delta)(x+1/(x\delta))}\Bigr),
\]
and
\[
N_+\le\int_0^\infty\Bigl(\frac x{1+x^2}-\frac{\delta x}{1+\delta x^2}\Bigr)\dd x+\int_0^\infty e^{-r(1-\delta)(x+1/(x\delta))}\frac{\dd x}x=\frac12\log\frac1\delta+2K_0\bigl(2(1-\delta)r/\sqrt\delta\bigr),
\]
which is $L+2K_0(2-2r^2)$.

\emph{$e_3$.} By \eqref{eq:H} and Lemma~\ref{lem:J}(b), with \eqref{eq:bessel} at $\nu=\pm1$,
\begin{equation}\label{eq:e3id}
\int_0^\infty\Bigl(1+\frac1{x^2}\Bigr)\partial_rJ\dd x=4r\int_0^\infty u^2K_1(2ru)\dd u=\frac1{r^2},
\end{equation}
\[
I_1:=\int_0^\infty\frac{\partial_rJ}{x^2}\dd x=4r\int_0^\infty\frac{u^2w}{1+w}K_1(2ru)\dd u.
\]
Hence $D_+'=2I_1-1/r^2$ and $e_3=4r^2I_1>0$. With $\frac{w}{1+w}\le\min(\frac12,u^{-2})$, splitting at $u=\sqrt2$ and using $u^2K_1(2ru)\le u/(2r)$ on the first piece, $I_1\le1+2K_0(2\sqrt2\,r)$.

\emph{$e_4$.} By \eqref{eq:bessel} at $\nu=0$, $M:=\int_0^\infty\partial_rJ\,x^{-1}\dd x=4r\int_0^\infty\frac{u^2\sqrt w}{1+w}K_0(2ru)\dd u$, and $u^2\sqrt w/(1+w)=u-2uw/(1+w)$. Since $4r\int_0^\infty uK_0(2ru)\dd u=1/r$, $e_4=2-2rM=16r^2\int_0^\infty\frac{uw}{1+w}K_0(2ru)\dd u>0$.
\end{proof}

\begin{proposition}\label{prop:small}
$H_+>0$ for $0<r\le0.01$.
\end{proposition}

\begin{proof}
Since $H=N'D-ND'=4H_+$, expanding gives exactly
\[
\frac{r^2H}4=(L-\gamma-1)+\frac{2e_2+2e_4+e_1e_4-2e_1-e_3N}4\ \ge\ (L-\gamma-1)-\frac{e_1}2-\frac{e_3N}4,
\]
by Lemma~\ref{lem:e}. For $r\le0.01$, $K_0(y)\le K_0(1)+\log(1/y)$ with $K_0(1)<0.4211$, and $4K_0(2-2r^2)<0.456$, so Lemma~\ref{lem:e} gives $e_1/2\le4r^2(2L+0.456)$, $e_3\le8r^2L$ and $N\le2L+0.456$. Hence
\[
\frac{e_1}2+\frac{e_3N}4\le e^{-2L}(2L+0.456)(2L+4),
\]
which is decreasing for $L\ge\log100$ and at most $0.0128$ there, while $L-\gamma-1\ge3.027$.
\end{proof}

The lower bound in Theorem~\ref{thm:main} for small $r$ is $G=N/D\ge r(2L-2\gamma)/(2+e_1)$, and the upper bound is $G\le rN/2$; both are from Lemma~\ref{lem:e}.

\section{Large rates}\label{sec:large}

\begin{lemma}\label{lem:large}
For $x,r>0$:
\begin{enumerate}[label=(\alph*)]
\item $0\le N_+\le\bar N:=\frac\pi{4r}+K_1(r)$ and $D_+\ge\frac\pi2$;
\item $\partial_rJ\le\frac{16}{r^2x^3}+xe^{-rx/2}$, so $\int_X^\infty\partial_rJ\dd x\le\frac8{r^2X^2}+e^{-rX/2}\bigl(\frac{2X}r+\frac4{r^2}\bigr)$;
\item $\int_{1/2}^1\partial_rJ\dd x\ge0.016\,(1-Q)/r^2$, where $Q=e^{-a}(1+a+a^2/2)$, $a=5r/2$.
\end{enumerate}
\end{lemma}

\begin{proof}
(a) The integrand of $N_+$ is $\frac1{xb}(x^2-bJ)\ge0$ by Lemma~\ref{lem:J}(a). Write
\[
\frac{x^2}b-J=rx\int_0^1e^{-rsb/x}\Bigl(1-e^{-rs^2/(x(1-s))}\Bigr)\dd s+rx\int_1^\infty e^{-rsb/x}\dd s .
\]
On $s\le\frac12$ use $1-e^{-c}\le c\le2rs^2/x$; the rest is at most $rx\int_{1/2}^\infty e^{-rsb/x}\dd s=\frac{x^2}be^{-rb/(2x)}$. So $x^2-bJ\le\frac{4x^3}{rb^2}+x^2e^{-rb/(2x)}$. Dividing by $xb$, using $x/b\le\frac12$, $\int_0^\infty x^2b^{-3}\dd x=\pi/16$ and $\int_0^\infty e^{-(r/2)(x+1/x)}\dd x=2K_1(r)$ gives the bound on $N_+$. Since $J\le x^2/b$, $1-J\ge1/b$, so $D_+\ge\int_0^\infty b^{-1}\dd x=\pi/2$.

(b) With $t=s/(1-s)$, $\varphi=\frac{xt}{1+t}+\frac tx$ and $\partial_rJ=r\int_0^\infty\frac{t^2}{(1+t)^2}e^{-r\varphi}\dd t$. For $t\le1$, $\frac{xt}{1+t}\ge\frac{xt}2$; for $t\ge1$, $\frac{xt}{1+t}\ge\frac x2$ and $\frac{t^2}{(1+t)^2}\le1$. So $\partial_rJ\le r\int_0^\infty t^2e^{-rxt/2}\dd t+re^{-rx/2}\int_0^\infty e^{-rt/x}\dd t$.

(c) For $x\in[\frac12,1]$ and $t\in[0,1]$, $\varphi\le t(x+1/x)\le\frac52t$ and $\frac{t^2}{(1+t)^2}\ge\frac{t^2}4$, so $\partial_rJ\ge\frac r4\int_0^1t^2e^{-5rt/2}\dd t=\frac{4(1-Q)}{125\,r^2}$.
\end{proof}

\begin{proposition}\label{prop:large}
$H_+<0$ for $r\ge5$.
\end{proposition}

\begin{proof}
By \eqref{eq:H},
\[
-H_+=\int_0^\infty\partial_rJ\,\Bigl[\frac{D_+}x+\frac{N_+}{x^2}-N_+\Bigr]\dd x .
\]
The bracket is at least $D_+/x-N_+$, so it is negative only for $x>D_+/N_+\ge X:=(\pi/2)/\bar N$, and there it is at least $-\bar N$. On $[\frac12,1]$ it is at least $\frac\pi2-\bar N$. For $r\ge5$, $\bar N<0.17$ and $X>9$, so by Lemma~\ref{lem:large}
\[
-r^2H_+\ \ge\ 0.016\,(1-Q)\Bigl(\frac\pi2-\bar N\Bigr)-\bar N\Bigl(\frac8{X^2}+e^{-rX/2}(2rX+4)\Bigr)=:\ell(r).
\]
As $r$ grows, $\bar N$ and $Q$ decrease, $X$ and $rX=\frac\pi2\,r^2/(\frac\pi4+rK_1(r))$ increase, and $v\mapsto e^{-v/2}(2v+4)$ is decreasing; so $\ell$ is increasing. Finally $\ell(5)>0.0089$.
\end{proof}

With Lemma~\ref{lem:large}(a), $G=N_+/D_+\le2\bar N/\pi=\frac1{2r}+\frac2\pi K_1(r)$.

\section{The middle range}\label{sec:middle}

On $[0.01,100]$ the sign of $G'$ is certified with ball arithmetic in Arb \cite{Arb} through python-flint, at 64 bits of precision. These computations use the formula \eqref{eq:G0} on the whole line, without the symmetry, with $f_0$ written as an integral over $u\in(0,\infty)$ for $x>0$ and over a finite $u$-interval for $x<0$, so they also check the reduction of Section~\ref{sec:reduction}.

For an interval $R=[r_{\rm lo},r_{\rm hi}]$ the parameter enters as a ball, so every enclosure holds for all $r\in R$ at once. The $x$-integrals of $N$, $D$ and of their first and second $r$-derivatives are split into three parts.
\begin{itemize}
\item A strip $|x|\le\varepsilon$ with $\varepsilon=10^{-5}\min(r_{\rm lo},1)$. The integrand is smooth but not analytic at $x=0$, since the branch $x<0$ carries $e^{-r(1+x^2)/|x|}$; it is bounded directly (for instance $0\le d\le 2(1+x^2)x^3/r$ for $x>0$, where $d=(1+x^2)f_0-1$).
\item Tails $|x|\ge X$ with $X=10^4/\min(r_{\rm lo},1)$ (and $10^5/\min(r_{\rm lo},1)$ for the derivatives), bounded by inequalities of the type of Lemma~\ref{lem:large}(b).
\item Graded panels $[a,2a]$ on each side between $\varepsilon$ and $X$. On each panel the integral is enclosed by the $24$-point Gauss--Legendre rule plus the remainder bound $\frac{64}{15}\,h\,M\rho^{-48}/(\rho^2-1)$ for the Bernstein ellipse $E_\rho$ with $\rho=1.6$ \cite[Theorem 19.3]{Tr}, where $h$ is the half-length and $M$ bounds the integrand on the boundary of $E_\rho$ (maximum modulus), evaluated on $48$ boxes covering it. On these ellipses $|\arg x|<\pi/4$, which keeps the complex continuation of the integrand analytic; its value at a complex $x$ is itself a rigorous Arb integral in $u$, built from $\exp$, $\operatorname{expm1}$ and division only.
\end{itemize}
Every strip and tail bound is a sum of powers of $1/r$ and terms $e^{-cr}$, $c>0$, hence decreasing in $r$, and is evaluated at $r_{\rm lo}$.

These give enclosures of $N,D,N',D'$, hence of $G'=(N'D-ND')/D^2$, and of $H'=N''D-ND''$ (the terms $N'D'$ cancel). An interval on which the enclosure does not decide a sign is bisected. The results are in Table~\ref{tab:cert}.

\begin{table}[h]
\centering
\begin{tabular}{@{}lp{0.42\textwidth}l@{}}
\toprule
range of $r$ & certified & intervals \\
\midrule
$[0.01,0.215]$ & $G'>0$ & 47, ratio $1.08$, bisected \\
$[0.215,0.25]$ & $H'<0$; $G'>0$ on $[0.215,0.226]$ and $G'<0$ on $[0.235,0.25]$ & 35, width $0.001$ \\
$[0.25,100]$ & $G'<0$ & 83, ratio $1.08$, bisected \\
\bottomrule
\end{tabular}
\caption{Certified signs. Each range is a contiguous chain of decided intervals.}\label{tab:cert}
\end{table}

On $[0.215,0.25]$, $H$ is strictly decreasing, so it vanishes at most once there. It is positive at $0.226$ and negative at $0.235$, so $G'$ has exactly one zero $r_*$ in $[0.215,0.25]$, and $r_*\in(0.226,0.235)$.

As an independent check of the code, fifteen decided intervals were recomputed with ordinary adaptive quadrature at both endpoints and the midpoint; the values of $G$ and the central differences of $G$ fell inside the stored enclosures in all $45$ cases.

\begin{proof}[Proof of Theorem~\ref{thm:main}]
By Proposition~\ref{prop:small} and Table~\ref{tab:cert}, $G'>0$ on $(0,r_*)$; by Table~\ref{tab:cert} and Proposition~\ref{prop:large}, $G'<0$ on $(r_*,\infty)$. The bounds on $G$ are those at the ends of Sections~\ref{sec:small} and~\ref{sec:large}.
\end{proof}

\section*{Acknowledgements and funding}
The research presented in this paper was supported by the European Research
Council (ERC) under the European Union's Horizon Europe research and
innovation programme (grant agreement No.~101041711), by the Simons
Foundation as part of the Collaboration on the Mathematical and Scientific
Foundations of Deep Learning, by Heights Labs, by Convex Nexus Capital, by
the Israel Science Foundation (grant number 2258/19), and by the Israel
Science Foundation (ISF Grant 4101/25). Computational resources were
provided in part by Google Cloud through its education credits program
(credit 530550103).

\section*{Declaration of competing interest}
The author holds equity in Heights Labs and in Convex Nexus Capital and
has a management role at Convex Nexus Capital. These entities are included
in the funding declaration above.

\section*{Declaration of generative AI and AI-assisted technologies}
The paper comes out of Math Brain, an ongoing project on the use of AI systems for research mathematics. In November 2026 Claude Opus~5.5 (Anthropic), working in Claude Code, found the proofs of Sections~\ref{sec:reduction}--\ref{sec:large}, wrote and ran the interval computations of Section~\ref{sec:middle} and the simulation of Remark~\ref{rem:scale}, and wrote the paper. The author directed the work and is responsible for the content of the paper and for decisions concerning its dissemination.

\section*{Data availability}
No data were used. The code for the interval computations, the independent checks and the simulation, with their outputs, is available at\\ \url{https://github.com/yspennstate/malrieu-switched-lyapunov-unimodal}.

\begin{thebibliography}{9}

\bibitem{Arb}
F.~Johansson, Arb: efficient arbitrary-precision midpoint-radius interval arithmetic, \emph{IEEE Trans. Comput.} \textbf{66} (2017), 1281--1292.

\bibitem{LMR}
S.~D.~Lawley, J.~C.~Mattingly and M.~C.~Reed, Sensitivity to switching rates in stochastically switched ODEs, \emph{Commun. Math. Sci.} \textbf{12} (2014), no.~7, 1343--1352.

\bibitem{La}
G.~Lagasquie, A note on simple randomly switched linear systems, arXiv:1612.01861, 2016.

\bibitem{Ma}
F.~Malrieu, Some simple but challenging Markov processes, \emph{Ann. Fac. Sci. Toulouse Math. (6)} \textbf{24} (2015), no.~4, 857--883. Also arXiv:1412.7516.

\bibitem{Tr}
L.~N.~Trefethen, \emph{Approximation Theory and Approximation Practice}, extended edition, SIAM, Philadelphia, 2019.

\end{thebibliography}

\end{document}
